\scanpage{106}
\mspage{52}
11) Digression : déployement d'un diviseur à croisements normaux.

Pour fixer les idées, je vais me placer dans le contexte des multiplicités analytiques lisses — mais on pourrait \unsure{aussi|voir}, en prenant des multiplicités analytiques lisses \uline{relatives} — ou en remplaçant l'hyp. analytique\unsure{s} par des hyp. topologiques ad-hoc — ou en passant au contexte schématique.....

Soit $X$ une multiplicité analytique lisse, et $D$ un diviseur sur $X$ à croisements normaux (hyp. purement locale en $(X,D)$, par définition). \deleted{Soit $G$ un}

Pour tout pt $x$ de $D$, l'ens des ``branches'' de $D$ passant par $x$ est un ensemble fini. $I(x)$, Par définition, les \uline{germes} \inserted{en $x$} \uline{de sous-multiplicités} \uuline{attaché\unsure{s|es}} \uuline{à $D$} sont les intersections d'un nb fini \deleted{de} parmi les $D_{x,i}\in I(x)$.
\margin{il vaut mieux prendre un $x\in X$, et poser $I(x)=\phi$ si $x\notin D$.}
Ces \unsure{germes} forment un ensemble ordonné \inserted{$\Delta_x\unsure{=(|(=}\Delta_x(X\unsure{,D}))$} antiisomorphe \unsure{à|:} $\mathcal P(I(x))$, \deleted{Si \unsure{$J$ est}}\inserted{la codimension} \deleted{\unsure{une|avec|ceux}} de $D_{x,J}$ $(J\subset I(x))$ étant donnée par
\[
(149)\qquad \operatorname{Codim}(D_{x,J},X_x)=\operatorname{card}J
\]
\margin{105}
